
\sourcenote{Genre: core rule (tier 1). Setting (valuations, partial
evaluation, elementwise-under-broadcast atom semantics) as fixed in the
\rulename{EqElim} section.}

\textbf{Claim} (\texttt{sound\_eqRefl}).\enspace The payload supplies a
term $t$ and two checked typing conditions: $\bshape_H(t)$ is
\emph{defined} (some shape $s$), and every free wire of $t$ is typed at a
\emph{finite} (non-extended) dtype. Then $\Phi \vdash t = t$ is
semantically valid for an \emph{arbitrary} premise $\Phi$ --- the rule is
premise-free, and $\Phi$ rides along unread (a fact true in the empty
context is true under any premise). Indeed the proof below never reads
$\Phi$, the admissibility of $\sigma$, or the premise-satisfaction
hypothesis at all: the conclusion holds at \emph{every} valuation.

That anything must be checked here at all is the point of the section.
Equality atoms are evaluation-level facts (\rulename{EqElim} setting), so
``$t = t$'' is not free: it asserts that $t$ \emph{evaluates} --- and
evaluation is partial.

\begin{blemma}[totality under the typing pack]\label{lem:eqrefl-total}
If $\bshape_H(t) = s$ and every free wire of $t$ is finitely typed, then
at every valuation $\sigma$ the evaluation $\mathrm{ev}_\sigma(t)$ is
defined, with shape $s$.
\end{blemma}
\begin{proof}
Structural induction on $t$. A wire read succeeds because the wire's
declared dtype is finite: the valuation's entries lie in the dtype's
value set, which for a finite dtype contains no $\pm\infty$, so the
whole-tensor rational read succeeds, at the declared shape. A constant
evaluates to its rank-0 tensor. For a compound term, the induction
hypotheses give defined operands whose shapes are, by the definition of
$\bshape$, either equal or joinable through rank-0 broadcast --- exactly
the domain on which the elementwise operations are total; the result
lands at the join. A projection's rank-1-and-in-range gates are part of
$\bshape$'s definedness, so the projected read is in range.
(\texttt{Term.evalQ\_defined\_of\_bshape\_fin}.)
\end{proof}

\begin{blemma}[reflexivity lifts]\label{lem:eqrefl-lift}
For any tensor $A$, the elementwise-under-broadcast equality relation
holds between $A$ and $A$.
\end{blemma}
\begin{proof}
The shapes are equal, so the relation reads the entry list zipped against
itself: every pair on the diagonal is $(x, x)$.
(\texttt{QTensor.rel2\_refl}.)
\end{proof}

\begin{bthm}[soundness of \rulename{EqRefl}]\label{thm:eqrefl-sound}
Under the payload's typing conditions, $\Phi \vdash t = t$ is
semantically valid for every $\Phi$.
\end{bthm}
\begin{proof}
Fix $\sigma$. Lemma~\ref{lem:eqrefl-total} gives
$\mathrm{ev}_\sigma(t) = A$ for some tensor $A$;
Lemma~\ref{lem:eqrefl-lift} closes the atom with both sides read as $A$.
\end{proof}

\begin{bremark}[why each typing condition is essential]
\label{rem:eqrefl-cex}
Drop the finite-wires condition: let $w$ be a wire of an
\emph{extended} dtype (its value set admits $+\infty$), and let $\sigma$
assign it a tensor with an infinite entry --- admissible, since nothing
constrains a bare wire. The rational read of $w$ fails, $t := w$ does not
evaluate, and the atom $w = w$ is \emph{false} at $\sigma$; with
$\Phi := \top$ the sequent $\top \vdash w = w$ is refuted. Drop
shape-definedness: take $t := w_1 + w_2$ with $w_1$ of shape $[2]$ and
$w_2$ of shape $[3]$, both finite --- every wire reads, but the
elementwise sum refuses the incompatible shapes, $t$ never evaluates, and
$\vdash t = t$ is again false everywhere. In both cases the un-gated rule
would assert an invalid sequent under a satisfiable (indeed empty)
premise --- in the first case the conclusion fails at the exhibited
valuation (a finite assignment would satisfy it); in the second it is
unsatisfiable outright. The typing pack is exactly the semantic price of
equality-by-evaluation; it is the same field family that repairs
\rulename{Norm}, there compounded with canonical-form equality.
\end{bremark}

\begin{bremark}[taxonomy note]\label{rem:eqrefl-tax}
Wherever \rulename{Norm} applies, \rulename{EqRefl} is its diagonal
($t_1 = t_2 = t$, canonical forms trivially equal). It is nonetheless
kept primitive: the development's frozen minimal fragment (\texttt{v0} in
the sources) excludes the polynomial engine outright, and the equational
kit must stand without it.
\end{bremark}

\begin{anchors}
Lemma \ref{lem:eqrefl-total} & \texttt{Term.evalQ\_defined\_of\_bshape\_fin}; payload fields \texttt{tyDef}, \texttt{fin}\\
finite-wire reads succeed & \texttt{Hypernet.Valuation.wireQ\_isSome}\\
Lemma \ref{lem:eqrefl-lift} & \texttt{QTensor.rel2\_refl}\\
Theorem \ref{thm:eqrefl-sound} & \texttt{sound\_eqRefl}; wired as core \texttt{RuleData.eqRefl} into \texttt{ProofObject.sound}\\
typing predicates & \texttt{Term.bshape}, \texttt{Term.finWires} (Core/Catalogue)\\
\end{anchors}
